\subsection{问题四分析}
问题三揭示的复杂轨道只是现象，问题四旨在判别该现象是否伴随初值敏感性。此处最需避免的误区是将有限时间窗内的不规则轨道误判为稳定混沌。若正李雅普诺夫指数仅在特定积分步长、时间窗口或初始条件下出现，则其更可能是条件性数值伪影，而非系统固有的动力学阈值。

因此，除基准参数扫描外，我们在保持方程组不变的前提下，调整积分时长、步长及局部初始条件，检验正李雅普诺夫指数区间的鲁棒性。若不同数值设置下的判别结果一致，则系统动力学复杂性增强的趋势具有可靠解释力；若结果随数值设置显著波动，则仅能得出"系统复杂波动程度增强"的保守结论。

